\documentclass[11pt, a4paper]{article}
\usepackage[utf8]{inputenc}
\usepackage[ngerman]{babel}
\usepackage[T1]{fontenc}
\usepackage{geometry}
\geometry{a4paper, left=25mm, right=25mm, top=25mm, bottom=25mm}
\usepackage{helvet}
\renewcommand{\familydefault}{\sfdefault}
\usepackage{enumitem}
\usepackage{titlesec}
\usepackage{hyperref}
\usepackage{xcolor}
\usepackage{amssymb}
\usepackage{array}
\usepackage{booktabs}

\titleformat{\section}{\normalsize\bfseries}{}{0em}{}[\vspace{-0.3em}]
\titleformat{\subsection}{\small\bfseries}{}{0em}{}

\setlength{\parindent}{0pt}
\setlength{\parskip}{0.5em}

\definecolor{linkcolor}{RGB}{0,70,140}
\hypersetup{colorlinks=true, linkcolor=linkcolor, urlcolor=linkcolor}

\begin{document}

% ============================================================
% HEADER
% ============================================================
\begin{center}
    {\large\bfseries Kurzexpos\'e zur Bachelorarbeit}\\[0.5em]
    {\normalsize\bfseries 3D Gaussian Splatting f\"ur Scan-to-Building Information Modeling (BIM):\\
        Segmentierung und Centerline-Extraktion linienf\"ormiger Infrastruktur}\\[0.8em]
    \small
    \begin{tabular}{ll}
        \textbf{Studiengang:} & Geovisualisierung (B.Eng.) \\
        \textbf{Datum:}       & \today                     \\
    \end{tabular}
\end{center}

\vspace{-0.5em}
\noindent\rule{\textwidth}{0.4pt}

% ============================================================
\section{1. Problemstellung und Motivation}

Bei der Verlegung von Gleichstrom-Erdkabeln (Direct Current, DC) ist die laufende Dokumentation der Trassenlage ein zentraler Bestandteil der Baukoordination. Ein Praxisbeispiel aus internen Dokumenten des \"Ubertragungsnetzbetreibers TenneT fordert hierbei eine Genauigkeit von $\pm 10$\,cm von der Kabelachse in Lage und H\"ohe, sowie $5$\,cm f\"ur Kontrollpunkte. Die heutige Praxis -- manuelle Vermessung mit Totalstation oder RTK-GNSS -- ist zeitaufw\"andig, erfordert Fachpersonal und liefert nur diskrete Punkte statt fl\"achendeckender 3D-Modelle. Gleichzeitig zeigt die Forschung, dass \textbf{3D Gaussian Splatting} (3DGS) als neue Methode der bildbasierten 3D-Rekonstruktion klassische Photogrammetrie in Rechenzeit und Detailtreue \"ubertreffen kann (Kerbl et al., SIGGRAPH 2023).

\textbf{Forschungsl\"ucke:} W\"ahrend 3DGS f\"ur visuelle Qualit\"at optimiert ist, fehlt eine systematische Untersuchung, ob sich aus 3DGS-Modellen \textbf{geometrisch verwertbare Achslinien} f\"ur r\"ohrenf\"ormige Infrastruktur ableiten lassen -- und wie genau diese im Vergleich zu geod\"atischen Referenzdaten sind.

% ============================================================
\section{2. Forschungsfragen}

\begin{enumerate}[noitemsep, leftmargin=1.5em]
    \item L\"asst sich aus einem 3DGS-Modell \"uber Mesh-Extraktion (SuGaR) und Centerline-Algorithmen (DGtal) eine geometrisch verwertbare 3D-Achslinie f\"ur r\"ohrenf\"ormige Infrastruktur ableiten?
    \item Welche der drei untersuchten Segmentierungsmethoden -- Farbfilterung (Baseline), FlashSplat (trainingsfrei), Gaussian Grouping (trainingsbasiert) -- liefert die robusteste Objektisolation bei geringstem Aufwand?
    \item Wie genau ist die extrahierte Centerline im Vergleich zur Referenz eines terrestrischen Laserscanners (TLS) gemessen im Root Mean Square Error (RMSE)?
    \item \textbf{Praxis-Problem:} Wie l\"asst sich der systematische H\"ohenfehler der Centerline algorithmisch korrigieren (Z-Shift), der durch die unvollst\"andige optische Erfassung des Kabels im Graben (Halbzylinder-Geometrie) entsteht?
    \item Ab welchem Objektdurchmesser scheitert die Pipeline?
\end{enumerate}

% ============================================================
\section{3. Methodischer Ansatz}

Die Arbeit entwickelt eine durchg\"angige \textbf{Scan-to-BIM-Pipeline}, die von der Bildaufnahme bis zur georeferenzierten 3D-Achslinie im GIS reicht. Als alternativer Scope wird zudem die automatische Georeferenzierung \"uber eindeutig GNSS-eingemessene Passpunkte betrachtet. Alle Bausteine basieren auf ver\"offentlichter, quelloffener Forschungssoftware:

\begin{center}
    \small
    \begin{tabular}{c p{4.5cm} p{5.5cm} p{4cm}}
        \toprule
        \textbf{\#} & \textbf{Schritt}                                 & \textbf{Tool / Methode}                & \textbf{Referenz}            \\
        \midrule
        1           & Bilderfassung                                    & Kamera (Testfeld Campus)               & --                           \\
        2           & Structure-from-Motion                            & COLMAP                                 & Sch\"onberger (CVPR 2016)    \\
        3           & 3DGS-Training                                    & gsplat                                 & Kerbl et al. (SIGGRAPH 2023) \\
        4           & Georeferenzierung                                & Helmert 7P (eigenes Skript)            & Klassische Geod\"asie        \\
        5           & Segmentierung                                    & \textit{Vergleichsstudie (3 Methoden)} & s. Abschnitt 4               \\
        6           & Mesh-Extraktion                                  & SuGaR                                  & Gu\'edon (CVPR 2024)         \\
        7           & Centerline-Extraktion                            & DGtal / Python                         & Kerautret (DGCI 2016)        \\
        8           & Geometrische H\"ohenkorrektur (Z-Shift)          & Python (Radius-Offset)                 & L\"osung Halbzylinder-Problem \\
        9           & Geoinformationssystem (GIS)-Import \& Evaluation & ArcGIS Pro / QGIS                      & --                           \\
        \bottomrule
    \end{tabular}
\end{center}

\textbf{Kernbeitrag (Der Automatisierungs-Vorteil):} Die Arbeit entwickelt eine durchg\"angige, hochautomatisierte \textbf{Scan-to-BIM-Pipeline}. Der gro\ss e Vorteil dieses Ansatzes liegt in der Verkettung: Aus einem einfachen Drohnenflug plus wenigen Passpunkten (Ground Control Points, GCPs) wird vollautomatisch ein segmentiertes, georeferenziertes 3D-Mesh erzeugt. Da dieses Mesh nicht via klassischer Photogrammetrie (Multi-View Stereo, MVS), sondern topologisch sauber aus den Oberfl\"achen der 3D Gaussians (SuGaR) abgeleitet wird, ist es geometrisch glatt genug, um daraus \textbf{direkt eine vektorisierte, GIS-f\"ahige 3D-Achslinie (Centerline)} zu extrahieren. Der wissenschaftliche Mehrwert liegt im quantitativen \textbf{Segmentierungsvergleich} und der \textbf{Evaluation gegen TLS-Referenzdaten}. Die technische Umsetzung erfolgt dabei methodisch sauber \"uber ein zentrales Python-basiertes Orchestrierungs-Skript, das alle Sub-Tools steuert und ein automatisiertes Time-Tracking erm\"oglicht.

% ============================================================
\section{4. Segmentierungsvergleich (Kernbeitrag)}

F\"ur die Isolation der Kabel-Gaussians werden drei Methoden mit steigender Komplexit\"at verglichen. F\"ur einen ersten Prototypen wird zun\"achst nur die Farbfilterung verwendet, da diese am einfachsten umzusetzen ist. Bei Erfolg wird die Segmentierungsqualit\"at auf die weiteren Methoden gesteigert:

\begin{center}
    \small
    \begin{tabular}{p{3.5cm} p{6.5cm} c p{3cm}}
        \toprule
        \textbf{Methode}      & \textbf{Typ}                                      & \textbf{Training} & \textbf{Paper}   \\
        \midrule
        (A) Farbfilterung     & Baseline (Hue-Saturation-Value (HSV)-Schwellwert) & Nein              & Eigene Impl.     \\
        (B) FlashSplat        & Optimierung auf 2D-Annotationen                   & Nein              & Shen (ECCV 2024) \\
        (C) Gaussian Grouping & Identity Encoding via DEVA/SAM                    & Ja                & Ye (ECCV 2024)   \\
        \bottomrule
    \end{tabular}
\end{center}

Dieser Vergleich zeigt systematisch, ab welcher Komplexit\"atsstufe sich der Mehraufwand f\"ur die Kabelisolation in besserer Centerline-Genauigkeit auszahlt.

% ============================================================
\section{5. Evaluation und Referenzdaten}

\begin{itemize}[noitemsep, leftmargin=1.5em]
    \item \textbf{RTK-GNSS} (Hochschulger\"ate): Einmessung von 8--12 GCPs f\"ur die Georeferenzierung ($<$\,1\,cm Genauigkeit).
    \item \textbf{Terrestrischer Laserscanner} (Hochschulger\"at): Dichte Referenz-Punktwolke des Testfelds. Daraus wird eine \textbf{TLS-Referenz-Centerline} extrahiert, gegen die die 3DGS-Centerline quantitativ evaluiert wird (Linie-zu-Linie RMSE, nicht nur Punkt-zu-Punkt).
\end{itemize}

Evaluationsdimensionen: (D1)~Geometrischer RMSE vs. TLS, (D2)~Visuelle Qualit\"at -- Peak Signal-to-Noise Ratio (PSNR) / Structural Similarity Index Measure (SSIM), (D3)~Segmentierungsqualit\"at -- Intersection over Union (IoU), (D4)~Rechenzeit und Kosten, (D5)~Skalierbarkeit (Tile-\"Ubergangs-RMSE).

% ============================================================
\section{6. Testfeld und Skalierbarkeit}

Das Testfeld wird auf dem \textbf{Campus} aufgebaut (ca. 20--30\,m L\"ange). Es besteht jedoch der explizite Wunsch, stattdessen echte Testdaten auf einer realen TenneT-Baustelle zu erfassen (vorbehaltlich interner Kl\"arung). Die Kamera wird am Campus auf ca. \textbf{2\,m H\"ohe} gef\"uhrt; bei einer realen Drohnenh\"ohe von 10\,m ergibt sich ein \textbf{Skalierungsfaktor von 5$\times$}. Eine solche Skalierung ver\"andert jedoch aufgrund der gr\"o\ss eren Entfernung in der Realit\"at wichtige Kameraparameter, weshalb Pixel eher verschwimmen. Echte Daten von einer Baustelle sind daher einer skalierten Alternative zwingend vorzuziehen. F\"ur den simulierten Fall sind die Testobjekt-Durchmesser so gew\"ahlt, dass sie bei 2\,m die gleiche Pixelabdeckung (Ground Sampling Distance, GSD) erzeugen wie reale Kabel bei 10\,m Drohnenh\"ohe:

\begin{center}
    \small
    \begin{tabular}{p{3.8cm} p{2.5cm} p{3.5cm} p{2.5cm} p{2cm}}
        \toprule
        \textbf{Testobjekt}     & \textbf{$\varnothing$ Campus} & \textbf{Simuliert} & \textbf{$\varnothing$ real (10\,m)} & \textbf{Beschaffung} \\
        \midrule
        Gartenschlauch (dick)   & $\sim$40\,mm                  & DC-Erdkabel        & 20\,cm                              & Baumarkt             \\
        Gartenschlauch (d\"unn) & $\sim$16\,mm                  & Steuerkabel        & 8\,cm                               & Vorhanden            \\
        Paketschnur / Seil      & $\sim$10\,mm                  & Lichtwellenleiter  & 5\,cm                               & Vorhanden            \\
        Nylonschnur             & $\sim$3\,mm                   & Grenztest          & $<$\,2\,cm                          & Vorhanden            \\
        \bottomrule
    \end{tabular}
\end{center}

Diese Abstufung (40\,mm $\rightarrow$ 16\,mm $\rightarrow$ 10\,mm $\rightarrow$ 3\,mm) erm\"oglicht die systematische Untersuchung der Forschungsfrage~4 (Mindestdurchmesser). Das Campus-Testfeld vermeidet Datenschutzprobleme und erm\"oglicht sofortige Datenverf\"ugbarkeit.

F\"ur reale Trassen (100\,m--1\,km) wird ein \textbf{Tile-basiertes Processing} konzipiert: Die Szene wird in Abschnitte mit \"Uberlappung unterteilt, jeder Tile durchl\"auft die Pipeline separat, und die georeferenzierten Centerlines werden im GIS zusammengef\"uhrt.

% ============================================================
\section{7. Machbarkeit und Ressourcen}

\begin{center}
    \small
    \begin{tabular}{p{6.5cm} l p{5.5cm}}
        \toprule
        \textbf{Ressource}                               & \textbf{Status}                       & \textbf{Anmerkung}                \\
        \midrule
        Graphics Processing Unit (GPU, RTX 4000, 20\,GB) & \textit{Vorhanden}                    & Ausreichend f\"ur alle Tools      \\
        RTK-GNSS \& TLS                                  & \textit{Hochschule}                   & Rechtzeitig reservieren           \\
        3DGS Pipeline Tools                              & \textit{Open Source}                  & COLMAP, gsplat, SuGaR, FlashSplat \\
        Testfeld (Campus)                                & \textit{vermutlich schnell m\"oglich} & Genehmigungsbedarf ?              \\
        \bottomrule
    \end{tabular}
\end{center}

Alle Software-Bausteine sind quelloffen, voll publiziert (SIGGRAPH, CVPR, ECCV 2023/2024) und auf der vorhandenen Hardware lauff\"ahig. Die technische Realisierbarkeit ist durch Fallback-Optionen abgesichert (z.\,B. Fokus auf die trainingsfreie Segmentierung mit FlashSplat, falls Gaussian Grouping den VRAM \"ubersteigt).

% ============================================================
\section{8. Abgrenzung PA / BA}

\begin{itemize}[noitemsep, leftmargin=1.5em]
    \item \textbf{Projektarbeit} (8 ECTS, vorangehend): Testfeldaufbau, Datenerfassung, COLMAP + gsplat einrichten, visueller Vergleich MVS vs. 3DGS. \textit{Ergebnis:} Validierter Datensatz und lauff\"ahige Basispipeline.
    \item \textbf{Bachelorarbeit} (15 ECTS): Georeferenzierung, Segmentierungsvergleich, Mesh-Extraktion, Centerline, quantitative Evaluation gegen TLS. \textit{Ergebnis:} Vollst\"andige Scan-to-BIM-Pipeline mit RMSE-Bewertung.
\end{itemize}

% ============================================================
\section{9. Ausgew\"ahlte Literatur}

\begin{enumerate}[noitemsep, leftmargin=1.5em, label={[\arabic*]}]
    \item Kerbl et al.: \textit{3D Gaussian Splatting for Real-Time Radiance Field Rendering.} SIGGRAPH 2023.
    \item Gu\'edon, Lepetit: \textit{SuGaR: Surface-Aligned Gaussian Splatting for Efficient 3D Mesh Reconstruction.} CVPR 2024.
    \item Ye et al.: \textit{Gaussian Grouping: Segment and Edit Anything in 3D Scenes.} ECCV 2024.
    \item Shen et al.: \textit{FlashSplat: 2D to 3D Gaussian Splatting Segmentation Solved Optimally.} ECCV 2024.
    \item Kerautret et al.: \textit{3D Geometric Analysis of Tubular Objects based on Surface Normal Accumulation.} DGCI 2016.
    \item He et al.: \textit{A Survey on 3DGS Applications: Segmentation, Editing, and Generation.} ArXiv 2508.09977, 2025.
    \item Sch\"onberger, Frahm: \textit{Structure-from-Motion Revisited.} CVPR 2016.
\end{enumerate}

\end{document}
